% Chapter 1
% 
\chapter{Introduction} % Main chapter title
\label{chap:Chapter1} % For referencing the chapter elsewhere, use Chapter~\ref{Chapter1}


%-------------------------------------------------------------------------------
%---------
%
\section{Writing guidelines} 
\label{sec:chap1_guidelines} %For referencing this section elsewhere, use Section~\ref{sec:chap1_guidelines}

This chapter presents some recommendations to follow in the structuring and formatting of the dissertation research plan.

The purpose of writing a research plan is, in essence, to document the work developed and highlight its importance. The author must follow the best writing rules, that is, respect grammatical and spelling rules, and correctly organize the text in order to make the message to be conveyed clear. 

The language used must be rigorous, scientific and without colloquial expressions. If in Portuguese, the document must be written using the rules of the new Orthographic Agreement that took effect in January 2009. If in English, the document should follow the rules of the United Kingdom variant. 

Automatic spelling correction must be complemented by consulting grammars and glossaries.

\section{Document structure}

The structure of the dissertation is normally divided into three main parts: Introduction, State of the art and Research Plan.

The Introduction should contain the following subsections:
\begin{itemize}
    \item Background and Motivation (context, relevance and motivation)
    \item Problem Statement (the problem, gaps and definition of scope)
    \item Goal and Objectives (the main goal and its measurable objectives)
    \item Research Questions (the questions that will guide the research)
    \item Privacy and Ethics (see Section 1.4)
    \item Use of AI Tools (see Section 1.5)
    \item Document Structure (brief description of the following chapters)
\end{itemize}

The main Body of the document includes the survey of the state of the art carried out, normally placed in the chapter following the Introduction. The survey of the state of the art should not be too exhaustive, focusing on the specific topics addressed in the document. This state of the art is intended to be updated during the writing of the dissertation. The structure of this chapter should be discussed with the supervisor.

An optional chapter can be added, with initial work already performed, such as a proposed approach and/or preliminary results.

Finally, a Research Plan chapter needs to be provided, with information on the work to be performed during the 2\textsuperscript{nd} semester.

\begin{itemize}
    \item the methodology of the work to be developed (summary of the approach, methods, prototypes, experimentation, etc.);
    \item how the results are going to be evaluated to allow the research questions to be answered;
    \item timeline of activities for the second semester.
\end{itemize}

\subsection{Contents checklist}

Table~\ref{tab:checklist} provides a list to check the contents of the document. 

\begin{table}[H]
\caption{Contents checklist.}
\label{tab:checklist}
\centering
\begin{tabular}{l l l}
\toprule
\tabhead{Description} & \tabhead{Mandatory?} & \tabhead{Page numbers} \\
\midrule
First page (template first page)&Yes&None\\
Statement of Integrity&Yes&Roman\\
Portuguese abstract (Resumo)&Yes&Roman\\
English abstract &Yes&Roman\\
Table of Contents&Yes&Roman\\
List of Figures&Yes&Roman\\
List of Tables&Yes&Roman\\
Acronyms and Symbols&Yes&Roman\\
Introduction&Yes&Arabic\\
State of the art&Yes&Arabic\\
Preliminary Results&No&Arabic\\
Research Plan&Yes&Arabic\\
References&Yes&Arabic\\
Annexes and appendices&No&Arabic\\

\bottomrule\\
\end{tabular}
\end{table}


\section{Formatting}

The document can be written in Portuguese or English. The minimum number of pages is 30 and the maximum is 60. Complementary documentation may be added in the form of annexes, never exceeding 90 pages in total (including the main text of the dissertation). 

Please follow the margins and fonts  defined in this template. The font size should be 11pt. The document should be printed double sided.

Please note that the graphic appearance and formatting are not a substitute for well-organized writing and a clear presentation of ideas. 

Please refer to the annex~\ref{AppendixA} for details about this template, how to format the document and insert citations, figures, tables, equations and other elements.

\section{Privacy and ethics}
The document is required to provide in the introduction chapter a specific section with information on how privacy and ethics were considered in the analysis and development of the work. If not applicable, this should be justified.  Without this section, the document will not be evaluated.

\section{Use of AI tools}
The dissertation course uses the model of limited use with declaration of use and intentions~\parencite{oliveira2025}. AI tools can be used, under supervision, to support the development of the work, such as idea generation, topic organization, document structure, brainstorming and linguistic revision. All content needs to be revised, and the author needs to add a statement taking full responsibility for the work. 

Direct use of content generated from AI tools is not allowed, unless the content itself is intended to be part of some experimental analysis, and clearly identified as so. 

The document is required to provide in the introduction chapter a specific section with information on which AI tools were used, for what purpose, and in which sections of the document. If no AI tools were used, the author should declare it in this section. Without this section, the dissertation will not be evaluated.

The following tables provide an example of summary tables that can be included in such section. These are mock-up examples, do not consider them as guidelines or recommendations. 

\begin{table}[htbp]
\centering
\caption{AI tools used}
\label{tab:ai-tools}
\begin{tabularx}{\textwidth}{l l l}
\toprule
Application & Model & Version \\
\midrule
ChatGPT (via chatgpt.com) & GPT-4o & Several versions since 2026-01-01 \\
GitHub Copilot (via VS Code) & GPT-4o, Claude Haiku 4.5 & 1.134 \\
llama.cpp & Gemma 4 & E4B \\
\bottomrule
\end{tabularx}
\end{table}

\begin{table}[htbp]
\centering
\caption{Purpose of use}
\label{tab:ai-purpose}
\begin{tabularx}{\textwidth}{l X}
\toprule
Chapter/Section & Purpose \\
\midrule
Abstract, Structure &
Since the contents depend on, and can be inferred from, the rest of the document, models were used to generate the sections. However, all generated content has been critically reviewed to guarantee correctness. \\

Introduction &
Research questions were brainstormed with the model. The final version is the result of the author's critical analysis of the brainstorm. \\

State of the Art &
State of the art was brainstormed with models. Following an initial list of references identified by the author using the PRISMA method, models were used to screen and prioritise references, and identify related and missing works. \\

State of the Art &
The structure of the state of the art chapter was initially brainstormed with models, but was later evolved and changed during the writing. \\

All &
Models were used to propose grammar corrections and text simplification. \\

All &
Models were used to verify the correctness and completeness of bibliography and cross-references. \\
\bottomrule
\end{tabularx}
\end{table}

